%\documentclass[12pt,twoside]{article}
\documentclass[conference]{IEEEtran}  % this should work with your LaTeX installation; else download extra package (www.ctan.org/pkg/ieeetran) or remove IEEE usage below

%%%%%%% Fill this out:
\newcommand{\trtitle}{(Fine-tuned BERT) Encoder Comparison for Document Classification}
\newcommand{\titleshort}{KPIS EC4DC} % title for header:
\newcommand{\authorlastnames}{Ebert, Hamza} % alphabetical for seminars
\newcommand{\trcourse}{Knowledge Processing in Intelligent Systems: Practical Seminar}
\newcommand{\trgroup}{Knowledge Technology, WTM}
\newcommand{\truniversity}{Department of Informatics, University of Hamburg}
%%%%%%%%%%%%%%%%%%%%%%%%%%%%%%%%%%%%%%%%%%%%%%%%%%%%%%%%%%%%%
% Languages:

% If the thesis is written in English:
\usepackage[english]{babel} 						
\selectlanguage{english}

%%%%%%%%%%%%%%%%%%%%%%%%%%%%%%%%%%%%%%%%%%%%%%%%%%%%%%%%%%%%%
% Bind packages:

\usepackage{acronym}                    % Acronyms
\usepackage{algorithmic}				% Algorithms and Pseudocode
\usepackage{algorithm}				% Algorithms and Pseudocode
\usepackage{amsfonts}                   % AMS Math Packet (Fonts)
\usepackage{amsmath}                    % AMS Math Packet
\usepackage{amssymb}                    % Additional mathematical symbols
\usepackage{amsthm}
\usepackage{booktabs}                   % Nicer tables
%\usepackage[font=small,labelfont=bf]{caption} % Numbered captions for figures
\usepackage{color}                      % Enables defining of colors via \definecolor
\definecolor{uhhRed}{RGB}{254,0,0}		  % Official Uni Hamburg Red
\definecolor{uhhGrey}{RGB}{122,122,120} % Official Uni Hamburg Grey
\usepackage{fancybox}                   % Gleichungen einrahmen
\usepackage{fancyhdr}					% Packet for nicer headers
%\usepackage{fancyheadings}             % Nicer numbering of headlines

%\usepackage[outer=3.35cm]{geometry} 	  % Type area (size, margins...) !!!Release version
%\usepackage[outer=2.5cm]{geometry} 		% Type area (size, margins...) !!!Print version
%\usepackage{geometry} 					% Type area (size, margins...) !!!Proofread version
\usepackage{geometry} 	  % Type area (size, margins...) !!!Draft version
\geometry{a4paper,body={7.0in,9.1in}}

\usepackage{graphicx}                   % Inclusion of graphics
%\usepackage{latexsym}                  % Special symbols
\usepackage{longtable}				% Allow tables over several parges
\usepackage{listings}                   % Nicer source code listings
\usepackage{multicol}					% Content of a table over several columns
\usepackage{multirow}					% Content of a table over several rows
\usepackage{rotating}					% Alows to rotate text and objects
%\usepackage[hang]{subfigure}            % Allows to use multiple (partial) figures in a fig
%\usepackage[font=footnotesize,labelfont=rm]{subfig}	% Pictures in a floating environment
\usepackage{tabularx}					% Tables with fixed width but variable rows
\usepackage{url,xspace,boxedminipage}   % Accurate display of URLs

%%%%%%%%%%%%%%%%%%%%%%%%%%%%%%%%%%%%%%%%%%%%%%%%%%%%%%%%%%%%%
% Configurationen:

\hyphenation{whe-ther} 				% Manually use: "\-" in a word: Staats\-ver\-trag

%\lstloadlanguages{C}                   % Set the default language for listings
\DeclareGraphicsExtensions{.pdf,.svg,.jpg,.png,.eps} % first try pdf, then eps, png and jpg
\graphicspath{{./src/}} 				% Path to a folder where all pictures are located
\pagestyle{fancy} 					% Use nicer header and footer

% Redefine the environments for floating objects:
\setcounter{topnumber}{3}
\setcounter{bottomnumber}{2}
\setcounter{totalnumber}{4}
\renewcommand{\topfraction}{0.9} 			  %Standard: 0.7
\renewcommand{\bottomfraction}{0.5}		  %Standard: 0.3
\renewcommand{\textfraction}{0.1}		  	%Standard: 0.2
\renewcommand{\floatpagefraction}{0.8} 	%Standard: 0.5

% Tables with a nicer padding:
\renewcommand{\arraystretch}{1.2}

%%%%%%%%%%%%%%%%%%%%%%%%%%%%
% Additional 'theorem' and 'definition' blocks:
\theoremstyle{plain}
\newtheorem{theorem}{Theorem}[section]
%\newtheorem{theorem}{Satz}[section]		% Wenn in Deutsch geschrieben wird.
\newtheorem{axiom}{Axiom}[section] 	
%\newtheorem{axiom}{Fakt}[chapter]		% Wenn in Deutsch geschrieben wird.
%Usage:%\begin{axiom}[optional description]%Main part%\end{fakt}

\theoremstyle{definition}
\newtheorem{definition}{Definition}[section]

%Additional types of axioms:
\newtheorem{lemma}[axiom]{Lemma}
\newtheorem{observation}[axiom]{Observation}

%Additional types of definitions:
\theoremstyle{remark}
%\newtheorem{remark}[definition]{Bemerkung} % Wenn in Deutsch geschrieben wird.
\newtheorem{remark}[definition]{Remark} 

%%%%%%%%%%%%%%%%%%%%%%%%%%%%
% Provides TODOs within the margin:
\newcommand{\TODO}[1]{\marginpar{\emph{\small{{\bf TODO: } #1}}}}

%%%%%%%%%%%%%%%%%%%%%%%%%%%%
% Abbreviations and mathematical symbols
\newcommand{\modd}{\text{ mod }}
\newcommand{\RS}{\mathbb{R}}
\newcommand{\NS}{\mathbb{N}}
\newcommand{\ZS}{\mathbb{Z}}
\newcommand{\dnormal}{\mathit{N}}
\newcommand{\duniform}{\mathit{U}}

\newcommand{\erdos}{Erd\H{o}s}
\newcommand{\renyi}{-R\'{e}nyi}
\usepackage{graphicx}

% correct bad hyphenation here
\hyphenation{}

%%%%%%%%%%%%%%%%%%%%%%%%%%%%%%%%%%%%%%%%%%%%%%%%%%%%%%%%%%%%%
% Document:
\begin{document}

\title{\trtitle\\[1.5ex]
\large \trcourse\\[0.5ex]
\trgroup\\[0.5ex]
\truniversity}

\renewcommand{\headheight}{14.5pt}
\fancyhead[LO]{\slshape \authorlastnames}
\fancyhead[RO]{ \slshape \titleshort}

\author{\IEEEauthorblockN{Sabrina Ebert}
\IEEEauthorblockA{\textit{sabrina.ebert@studium.uni-hamburg.de}}
\and
\IEEEauthorblockN{Muhammad Hamza}
\IEEEauthorblockA{\textit{muhammad.hamza-2@studium.uni-hamburg.de}}
}

\maketitle

% Outline (assignment 3): 4 pages + appendix

\begin{abstract}
Encoder models like BERT are small and cheap compared to large language models and are often used for classifying texts.
LoRA makes fine-tuning them cheaper by training only a few small added matrices.
It comes with several settings, such as rank, learning rate and which layers it is added to, and it is not clear to us which of them matter most for a small encoder.
In this work we fine-tune BERT-base with LoRA for classifying sentences, short news texts and full documents (SST-2, AG News, 20 Newsgroups) and measure if and for which settings the classification performance changes, and by how much.
We change one setting at a time, run every setting with three seeds and compare accuracy and macro F1 to a frozen BERT baseline.
Our working hypothesis is that learning rate and LoRA placement change the performance more than the rank does.
If time allows, we also take a first look inside the fine-tuned models with explainable AI methods, e.g.\ word attributions, to explain the results.
\end{abstract}

\IEEEpeerreviewmaketitle

\section{Introduction and Motivation}
\label{sec:introduction}
\begin{itemize}
  \item encoders small and cheap vs.\ LLMs
  \item LoRA: many settings, little guidance
\end{itemize}

\section{Theoretical Background}
\label{sec:background}
\subsection{BERT}
\subsection{LoRA}
\subsection{Related Work}

\section{Experimental Setup}
\label{sec:setup}
\begin{itemize}
  \item Kaggle, SST-2, AG News, 20 Newsgroups
  \item baseline, LoRA study, other datasets
  \item libraries: transformers, peft, scikit-learn, maybe captum (XAI)
  \item 3 seeds per setting
\end{itemize}

\section{Results}
\label{sec:results}

\section{Discussion and Outlook}
\label{sec:discussion}

\section*{Acknowledgments and Use of Tools and AI}
We used Claude Code (Anthropic) to set up the LaTeX layout, to suggest some of the references and to draft the abstract from our milestones and improve its grammar and wording.
The topic, study design and working hypothesis are our own.
We checked all text and references and take full responsibility for the content.

\nocite{devlin2019bert,hu2022lora,sun2019finetune,dodge2020finetuning,socher2013sst,wang2019glue,zhang2015charcnn,lang1995newsweeder}

\clearpage
\bibliographystyle{plain}
\bibliography{bib}

\section*{Appendix: Code}

\end{document}
